% ============================================================
% MÓDULO: MI PRIMER EMPLEO
% ============================================================
\clearpage
\section*{MÓDULO: MI PRIMER EMPLEO}
\subsection*{Del currículum a la entrevista}

\begin{recuadro}
\textbf{Objetivo del módulo:} que cada estudiante egrese con su currículum armado, su perfil
de LinkedIn creado, cuentas activas en los principales portales de empleo y la práctica de
haber atravesado un proceso de selección simulado (postulación, evaluación psicotécnica y
entrevista).
\end{recuadro}

% ------------------------------------------------------------
\section*{1. EL CURRÍCULUM VITAE (CV)}

El \textbf{currículum} es un documento breve que resume quién sos, qué sabés hacer y cómo
contactarte. Su único objetivo es \textbf{conseguir una entrevista}. Un reclutador le dedica a
cada CV muy pocos segundos en la primera lectura, por eso tiene que ser claro, ordenado y
fácil de leer.

\subsubsection*{1.1. Estructura recomendada}
\begin{enumerate}
  \item \textbf{Datos personales y de contacto:} nombre y apellido, localidad (no hace falta
        la dirección completa), teléfono, correo electrónico profesional
        (\texttt{nombre.apellido@...}, nunca un apodo) y el enlace a tu perfil de LinkedIn.
  \item \textbf{Perfil o resumen (3 o 4 renglones):} quién sos y qué buscás. Ejemplo:
        \textit{``Técnico en Aviónica egresado de la E.E.S.T. N° 7 de Quilmes, con práctica en
        ensayo y reparación de equipos de comunicaciones y navegación. Busco mi primera
        experiencia en un taller aeronáutico.''}
  \item \textbf{Formación:} título técnico, escuela y año de egreso. Si estás cursando una
        carrera, indicá el año en curso.
  \item \textbf{Experiencia:} empleos, pasantías, Prácticas Profesionalizantes, trabajos
        informales relacionados (reparación de equipos, instalaciones). Para cada una: lugar,
        período y qué hiciste, empezando por un verbo (\textit{reparé, instalé, armé, verifiqué}).
  \item \textbf{Conocimientos técnicos:} instrumental (osciloscopio, multímetro, generador de
        señales), soldadura SMD, lectura de manuales ATA, software (Office, Odoo, CAD),
        programación.
  \item \textbf{Idiomas:} con nivel real (básico, intermedio, avanzado). En aviación el inglés
        técnico es muy valorado.
  \item \textbf{Proyectos y cursos:} el proyecto final de la escuela, cursos y certificaciones.
\end{enumerate}

\subsubsection*{1.2. Reglas de oro}
\begin{itemize}
  \item \textbf{Una página} para un primer empleo.
  \item \textbf{Sin datos innecesarios:} no hace falta DNI, estado civil, religión ni foto,
        salvo que el aviso la pida.
  \item \textbf{Sin errores de ortografía.} Pedile a alguien que lo lea antes de enviarlo.
  \item \textbf{Verdadero:} todo lo que figura puede ser preguntado en la entrevista.
  \item \textbf{Adaptado a cada aviso:} se ajustan el perfil y las palabras clave según el
        puesto.
  \item \textbf{En PDF}, con un nombre de archivo claro: \texttt{CV\_Apellido\_Nombre.pdf}.
\end{itemize}

% ------------------------------------------------------------
\section*{2. CÓMO PASAR LOS FILTROS ATS}

\begin{recuadro}
Un \textbf{ATS} (\textit{Applicant Tracking System}, sistema de seguimiento de candidatos) es
el software que usan las empresas y los portales para recibir, ordenar y filtrar los CV.
Antes de que una persona lea tu currículum, muchas veces lo lee un programa que busca
\textbf{palabras clave} del aviso y descarta los que no coinciden.
\end{recuadro}

\paragraph*{Cómo lograr que el ATS lea bien tu CV.}
\begin{itemize}
  \item \textbf{Usá las palabras del aviso.} Si el aviso pide ``mantenimiento de equipos de
        aviónica'' y ``lectura de diagramas eléctricos'', esas frases tienen que aparecer tal
        cual en tu CV (si es verdad que sabés hacerlo).
  \item \textbf{Formato simple:} una sola columna, títulos estándar (Experiencia, Formación,
        Conocimientos), letra común. Las tablas, columnas, íconos, gráficos de habilidades y
        textos dentro de imágenes suelen leerse mal o no leerse.
  \item \textbf{Archivo de texto real:} un PDF generado desde un procesador de textos, nunca
        un CV escaneado o una foto.
  \item \textbf{Siglas y nombre completo:} ``Taller Aeronáutico de Reparación (TAR)'', para que
        coincida cualquiera de las dos formas que use el aviso.
  \item \textbf{Fechas claras:} mes y año (03/2026 -- 12/2026).
  \item \textbf{Completá todos los campos} del formulario del portal, aunque repitan lo que ya
        dice el CV: el ATS filtra por esos campos.
\end{itemize}

\paragraph*{El uso de inteligencia artificial.}
Las herramientas de IA ayudan a corregir la redacción o adaptar el CV a un aviso, pero lo
que escriben hay que revisarlo: no deben inventar experiencia ni conocimientos, y un texto
genérico se nota.

% ------------------------------------------------------------
\section*{3. DÓNDE BUSCAR TRABAJO}

\begin{center}
\small
\begin{tabular}{|p{0.20\linewidth}|p{0.70\linewidth}|}
\hline
\textbf{Plataforma} & \textbf{Características} \\ \hline
LinkedIn & Red social profesional. Muchas empresas la usan como canal principal para
buscar personal. Sirve también para construir una red de contactos y una marca personal. \\ \hline
Computrabajo & El portal con mayor cantidad de avisos del país. Fuerte en perfiles
operativos, técnicos, atención al cliente y comercio. \\ \hline
Bumeran y ZonaJobs & Pertenecen al mismo grupo (Jobint). Fuertes en perfiles
administrativos, comerciales, técnicos y juniors; ZonaJobs tiene foco en Buenos Aires. \\ \hline
Portal Empleo (Nación) y oficinas municipales de empleo & Servicios públicos y gratuitos de
intermediación laboral, capacitación y programas para jóvenes. \\ \hline
Indeed y Google Empleos & Metabuscadores: reúnen avisos de muchos portales y permiten crear
alertas. \\ \hline
Sitios de las empresas & Aerolíneas, talleres y empresas del sector publican búsquedas en su
sección ``Trabaja con nosotros''. \\ \hline
\end{tabular}
\end{center}

\paragraph*{Consejos para usar los portales.}
\begin{itemize}
  \item Completá el perfil al 100\% y subí tu CV en PDF.
  \item Configurá \textbf{alertas} con palabras clave (``técnico aviónico'', ``técnico
        electrónico'', ``mantenimiento aeronáutico'', ``pañolero'').
  \item Postulate solo a avisos donde cumplas la mayoría de los requisitos.
  \item Llevá un \textbf{registro de postulaciones} (empresa, puesto, fecha, estado).
  \item \textbf{Cuidado con las estafas:} ninguna empresa seria cobra por postularse,
        capacitar o ``reservar'' un puesto, ni pide datos bancarios o claves antes de
        contratar.
\end{itemize}

% ------------------------------------------------------------
\section*{4. LINKEDIN}

LinkedIn es tu \textbf{currículum público} y tu red de contactos profesionales.

\subsubsection*{4.1. Un perfil completo}
\begin{itemize}
  \item \textbf{Foto profesional:} fondo neutro, buena luz, de frente. No una selfie de fiesta.
  \item \textbf{Titular:} lo que sos y lo que hacés, no solo ``Estudiante''. Ejemplo:
        \textit{``Técnico en Aviónica | Mantenimiento de equipos de comunicaciones y navegación
        | Soldadura SMD''}.
  \item \textbf{Acerca de:} un resumen de 4 a 6 renglones en primera persona.
  \item \textbf{Experiencia y formación:} igual que en el CV, con la escuela vinculada.
  \item \textbf{Aptitudes:} entre 10 y 20 habilidades concretas.
  \item \textbf{Destacados:} fotos o enlaces del proyecto final (sin datos reservados de
        empresas).
  \item \textbf{URL personalizada:} \texttt{linkedin.com/in/nombre-apellido}.
  \item Activar la opción \textbf{``Open to work''} (disponible para trabajar), visible solo
        para reclutadores si se prefiere.
\end{itemize}

\subsubsection*{4.2. Cómo usarlo}
\begin{itemize}
  \item Conectar con compañeros, docentes, instructores de prácticas y profesionales del
        sector, siempre con un mensaje breve de presentación.
  \item Seguir empresas del sector (aerolíneas, talleres, fabricantes, aeropuertos).
  \item Publicar de vez en cuando avances del proyecto final, cursos terminados o una visita
        técnica.
  \item Pedir \textbf{recomendaciones} a docentes o tutores de prácticas.
\end{itemize}

% ------------------------------------------------------------
\section*{5. CÓMO APLICAR A UN TRABAJO}

\begin{enumerate}
  \item \textbf{Leer el aviso completo} y subrayar requisitos excluyentes y deseables.
  \item \textbf{Investigar la empresa:} qué hace, dónde está, qué tipo de aeronaves o equipos
        maneja.
  \item \textbf{Adaptar el CV} con las palabras clave del aviso.
  \item Escribir una \textbf{carta o mensaje de presentación} breve: por qué te interesa ese
        puesto y qué podés aportar (tres párrafos cortos como máximo).
  \item \textbf{Postularse} completando todos los campos del formulario.
  \item \textbf{Registrar} la postulación y hacer un seguimiento amable si pasan dos o tres
        semanas sin respuesta (cuando hay un contacto disponible).
\end{enumerate}

% ------------------------------------------------------------
\section*{6. EL PROCESO DE SELECCIÓN}

El área de Recursos Humanos (RR.\,HH.) sigue, en general, estas etapas:

\begin{enumerate}
  \item \textbf{Definición del puesto:} el área que necesita personal describe tareas,
        requisitos y condiciones (el \textit{perfil del puesto}).
  \item \textbf{Reclutamiento:} se publica la búsqueda (portales, LinkedIn, consultoras,
        referidos, escuelas técnicas).
  \item \textbf{Preselección:} se filtran los CV, muchas veces con un ATS, y a veces se hace
        una breve entrevista telefónica.
  \item \textbf{Evaluaciones:} pruebas técnicas, exámenes psicotécnicos, dinámicas grupales.
  \item \textbf{Entrevistas:} con RR.\,HH. y con el jefe del área (entrevista técnica).
  \item \textbf{Verificación de referencias y examen preocupacional.} En aviación puede
        requerirse además la Certificación Médica Aeronáutica y la habilitación de seguridad
        aeroportuaria para trabajar en plataforma.
  \item \textbf{Oferta e incorporación:} propuesta formal, firma del contrato e inducción.
\end{enumerate}

\subsubsection*{6.1. Exámenes psicotécnicos}
Evalúan aptitudes, personalidad y forma de trabajar. Los más frecuentes son:
\begin{itemize}
  \item \textbf{Tests de aptitud:} razonamiento lógico, numérico, verbal, abstracto (series de
        figuras), atención y concentración. En puestos técnicos suelen incluir
        \textbf{razonamiento mecánico y espacial}.
  \item \textbf{Tests de personalidad:} cuestionarios sobre preferencias y conductas.
  \item \textbf{Tests proyectivos:} dibujos (persona bajo la lluvia, casa--árbol--persona) o
        completar frases.
  \item \textbf{Dinámicas grupales:} resolver un caso en equipo mientras se observa la
        comunicación, el liderazgo y la colaboración.
\end{itemize}
\textbf{Recomendaciones:} dormir bien la noche anterior, llegar puntual, leer bien cada
consigna, administrar el tiempo y responder con sinceridad los tests de personalidad: no hay
respuestas ``correctas'', y las respuestas forzadas suelen detectarse por incoherencias.

\subsubsection*{6.2. La entrevista}
\begin{itemize}
  \item \textbf{Antes:} investigar la empresa, releer el aviso y tu CV, preparar ejemplos
        concretos y tener preguntas para hacer.
  \item \textbf{Presentación:} ropa prolija acorde al puesto, puntualidad (o conexión probada si
        es virtual), celular en silencio.
  \item \textbf{Preguntas frecuentes:} ``Contame de vos'', ``¿Por qué este puesto?'',
        ``¿Cuáles son tus fortalezas y aspectos a mejorar?'', ``Contame una situación en la
        que resolviste un problema''.
  \item \textbf{Método STAR} para responder con ejemplos: \textbf{S}ituación, \textbf{T}area,
        \textbf{A}cción que realizaste y \textbf{R}esultado obtenido.
  \item \textbf{Entrevista técnica:} pueden pedir que identifiques componentes, leas un
        diagrama, uses instrumental o expliques el funcionamiento de un sistema. Si no sabés
        algo, decilo y explicá cómo lo averiguarías.
  \item \textbf{Después:} un breve mensaje de agradecimiento.
\end{itemize}

% ------------------------------------------------------------
\section*{7. NETWORKING ENTRE PARES}

El \textbf{networking} es construir y cuidar una red de contactos profesionales. Muchos
empleos se consiguen por \textbf{referidos}: alguien que conoce tu trabajo te recomienda.

\begin{itemize}
  \item \textbf{Tus compañeros son tu primera red:} dentro de pocos años van a estar trabajando
        en talleres, aerolíneas y empresas. Mantener el contacto y recomendarse mutuamente es
        una práctica habitual.
  \item \textbf{Docentes, tutores de prácticas y egresados} de la escuela: pueden avisarte de
        búsquedas y darte referencias.
  \item \textbf{Eventos del sector:} jornadas, ferias, visitas a talleres, cursos.
  \item \textbf{Dar antes de pedir:} compartir avisos, ayudar a un compañero con su CV,
        recomendar a quien hizo un buen trabajo.
  \item \textbf{Cuidar la reputación:} en un sector chico como el aeronáutico, la
        responsabilidad, la puntualidad y la ética se conocen y se comentan.
\end{itemize}

% ------------------------------------------------------------
\section*{8. ACTIVIDADES DEL MÓDULO}

\begin{enumerate}
  \item \textbf{Mi CV:} cada estudiante arma su CV en una página siguiendo la estructura del
        punto 1 y lo entrega en PDF.
  \item \textbf{Prueba ATS:} se elige un aviso real de un portal. Cada estudiante marca las
        palabras clave del aviso y adapta su CV. En parejas, se revisan los CV y se cuenta
        cuántas palabras clave coinciden.
  \item \textbf{Perfil de LinkedIn:} crear o mejorar el perfil, con titular, ``Acerca de'',
        aptitudes y URL personalizada. Conectarse con todo el curso.
  \item \textbf{Portales:} crear el perfil en Computrabajo y en Bumeran o ZonaJobs, y
        configurar dos alertas de búsqueda.
  \item \textbf{Simulacro de selección:} en grupos, un equipo arma un aviso y hace de RR.\,HH.
        (preselección y entrevista) y otro se postula. Después se invierten los roles. Se
        incluye una prueba técnica corta y una dinámica grupal.
  \item \textbf{Mapa de red:} cada estudiante arma una lista de 10 contactos profesionales
        (compañeros, docentes, tutores, empresas) y cómo podría ayudar y ser ayudado por cada
        uno.
\end{enumerate}
